\section{Related Work}
\label{sec:related}

This section draws only on entries marked VERIFIED in two literature
memos compiled for this paper
(\texttt{research/ncr\_separation\_grounding.md},
\texttt{research/ortho\_write\_grounding.md}, both dated 2026-07-16, live
web-search verification against arXiv/OpenReview, cross-checked against
this program's own independent prior sweep in
\texttt{NOVEL\_ARCH\_WATERFALL.md} \S1--\S5).

\textbf{Composition mechanism.} \citet{schlag2021fwm}'s Fast Weight
Memory is the single closest prior art on in-context-written,
composable fast weights: differentiable per-step writes into an
associative fast-weight matrix for compositional inference. Its reads are
recursive, LayerNorm-normalized, and fixed at a small hop count ($N_r{=}3$
in the original paper) -- approximate, not exact, and not a matrix-power
read; it is also the FWM-read baseline this paper's own Axis-A comparison
beats (Section~\ref{sec:win}). Log-Linear Attention
(\href{https://arxiv.org/abs/2506.04761}{arXiv:2506.04761}, ICLR 2026)
achieves logarithmic \emph{compute} complexity via a hierarchical,
Fenwick-tree-like summarization of the hidden state -- a training/inference
cost property of the write's structure, not a query-time repeated-squaring
read of one already-written operator matrix, and it targets standard
retrieval benchmarks rather than relational composition. HOLA
(\href{https://arxiv.org/abs/2607.02303}{arXiv:2607.02303}) augments a
decayed linear-attention state with a bounded, exact key--value cache of
high-surprise tokens, read via sharpened softmax attention: ``exact''
there means verbatim storage, not algebraic composition of operators.
\citet{marchetti2026sequential} study how a network's \emph{weights} come
to represent a finite group's irreducible representations across a
sequence, proving perfect in-weights learning needs hidden width
exponential in sequence length -- the in-weights complement to NCR's
in-context setting, not a competitor.
% evidence: research/ncr_separation_grounding.md Part 2

\textbf{State-tracking expressivity, correctly not our headline.}
\citet{grazzi2025unlocking} prove that restricting a linear RNN's
transition eigenvalues to $[0,1]$ (the Mamba/DeltaNet default) forecloses
parity and other group state-tracking tasks (parity itself being the
solvable $\mathbb{Z}_2$ case), and that extending the eigenvalue range
fixes it; \citet{peng2025rwkv7}'s RWKV-7 achieves
NC$^1$-hard state tracking with a non-diagonal, input-dependent transition.
Both already establish that matrix-valued recurrent state can escape the
$\mathrm{TC}^0$ ceiling transformers face \citep{merrill2023parallelism} --
via an $O(h)$ sequential rollout. NCR's claim is orthogonal: given an
already-written state, an $O(\log h)$ query-time read, with exactness as
its necessary empirical companion. We do not claim a new expressivity
result and would be incorrect to.
% evidence: research/ncr_separation_grounding.md Part 3

\textbf{The collapse-vs-hold genre on a harder task family, and what
``exact'' does and does not mean (cite-and-distinguish).} This paper's
own construction is a single cyclic group ($K$ entities on one
Hamiltonian cycle, Section~\ref{sec:background}) -- deliberately simpler
than, and explicitly not an instance of, the non-solvable-group word
problems ($S_5$) that a separate, harder literature already studies; we
do not conflate the two. That harder genre is nonetheless the closest
external precedent for the general \emph{shape} this paper's $K$-axis
also shows -- in-distribution success collapsing past some scale or
depth -- and is an established pattern since 2022, not a phenomenon this
paper discovers. \citet{liu2023shortcuts} originate it on $S_5$ itself
(ICLR 2023 oral; arXiv v1 Oct.\ 2022): low-depth transformers construct
an exact finite-state shortcut that is nonetheless brittle off the
training distribution. \citet{li2025track} locate two distinct
\emph{learned} collapse mechanisms (an associative-scan construction and
a parity-pruned one) for $S_3$/$S_5$ state tracking inside real language
models. Closest empirically: \citet{yau2026sequential}'s Prefix-Scannable
Models hold an $S_5$ ``cups and balls'' task to test length $180$, where
a transformer and a Mamba baseline both collapse -- the nearest published
precedent to a collapse-vs-hold separation on a non-solvable-group task,
though theirs is a learned, observed hold with no exactness argument,
unlike the by-construction read guarantee this program's mechanism
targets on its own, simpler task. \citet{siems2025deltaproduct}'s
DeltaProduct and the concurrent matrix-valued-state line
\citet{mishra2026m2rnn} sit in the same adjacent family; \citet{lee2026falsifier}
is further adjacent concurrent work, cited here for existence only, not
for any specific claim. One distinction is
load-bearing for this paper's own diagnosis
(Section~\ref{sec:diagnosis}) regardless of task family: DeltaProduct's
and RWKV-7's state-tracking results \citep{siems2025deltaproduct,
peng2025rwkv7} are \emph{expressivity} theorems -- a setting of weights
exists under which the recurrence eventually computes the correct group
product -- not a claim, enforced or measured, that an already-trained
model composes exactly at arbitrary query-chosen depth from a single
training regime; whether \emph{this program's own} trained write is
exact and well-conditioned, not whether an exact write could in
principle exist somewhere, is exactly what the Part~A/B ortho-write
verdict (Section~\ref{sec:pending}) measures. Finally, the $O(\log h)$
repeated-squaring read this paper builds on is not new mathematics: it
is the classical mechanism behind Barrington's theorem that bounded-width
branching programs capture $\mathrm{NC}^1$ via iterated, doubled group
composition \citep{barrington1989}; this paper studies whether a model
can be \emph{trained} to populate that read with a well-conditioned
operator, not the read's asymptotic complexity, which predates this
program by four decades and is unrelated to whether the composed group is
solvable.
% evidence: research/ncr_separation_grounding.md 2026-07-16 novelty-gate re-sweep;
% matrix-thinking/NCR_REAL_LM_DESIGN.md S N1 R4 (tiered must-cite list)

\textbf{Rank and readout.} The exact-continuous-recovery discipline of
Section~\ref{sec:background} follows \citet{nichani2025understanding}'s
rank-$m$-to-$\approx md$-associations argmax-capacity construction
\textbf{[NEEDS-FINAL-SPOT-CHECK, see Section~\ref{sec:background}]}.
% evidence: research/ncr_separation_grounding.md item 8

\textbf{Orthogonal parametrizations.} A line of unitary/orthogonal RNN
work \citep{arjovsky2016unitary,helfrich2018orthogonal} constrains the
recurrent transition matrix's eigenvalues to the unit circle; the Muon
optimizer (Jordan et al., 2024, code/blog release, no formal preprint;
validated at trillion-parameter scale by the Kimi K2 technical report,
\citealp{kimiteam2025k2}) orthogonalizes optimizer momentum updates via a
quintic Newton--Schulz iteration. Closest and most load-bearing to
disclose: \citet{nguyen2026muonssm}'s MuonSSM applies a single-iteration
Newton--Schulz projection to rank-1 fast-weight input injections for
general long-horizon SSM stability -- the differentiation from our
full-rank, 40-iteration, compositional-depth-motivated fix is given in
Section~\ref{sec:diagnosis}. The $(1.5,-0.5)$ cubic Newton--Schulz
coefficients we use trace to \citet{bjorck1971iterative}, not to Muon's
quintic form; the spectral pre-scaling step follows
\citet{higham1986computing}.
% evidence: research/ortho_write_grounding.md S1-S3

\textbf{This program's own, unrelated, older K-wall} (Task E
matrix-permutation composition, \emph{not} NCR) was rescued by training
budget alone; the distinction from this paper's K-wall, which budget alone
does not rescue, is disclosed in Section~\ref{sec:background}.
